% =============================================================================
% Chapter 10: Performance Engineering — ML Systems Lecture Slides (Vol II)
% =============================================================================
% IMAGES NEEDED (generate SVGs, convert to PDF):
%   - memory-wall.pdf          - roofline-model.pdf
%   - operator-fusion.pdf      - flashattention.pdf
%   - precision-engineering.pdf - graph-compilation.pdf
%   - optimization-playbook.pdf - diagnostic-flow.pdf
% =============================================================================
\documentclass[aspectratio=169, 12pt]{beamer}
\usepackage{../../assets/beamerthememlsys}

\mlsyssetup{
  volume       = {Volume II},
  chapter      = {Chapter 10},
  logo         = {../../assets/img/logo-mlsysbook.png},
  instlogo     = {../../assets/img/logo-harvard.png},
  chaptertitle = {Performance Engineering},
}

% --- Fonts ---
\usepackage{fontspec}
\setsansfont{Helvetica Neue}[
  BoldFont={Helvetica Neue Bold},
  ItalicFont={Helvetica Neue Italic},
  BoldItalicFont={Helvetica Neue Bold Italic},
]
\setmonofont{JetBrains Mono}[Scale=0.85]

% --- Packages ---
\usepackage{booktabs}
\usepackage{amsmath}

% --- Image paths ---
\graphicspath{
  {images/}
}

% --- Chapter-specific macros ---
\newcommand{\AI}{\ensuremath{I}}  % Arithmetic Intensity
\newcommand{\ridge}{\ensuremath{I_{\text{ridge}}}}

% --- Helper: safe image include ---
\newcommand{\safeimg}[2][width=\textwidth,keepaspectratio]{%
  \IfFileExists{images/#2}{\includegraphics[#1]{#2}}{%
    \IfFileExists{#2}{\includegraphics[#1]{#2}}{%
      \fbox{\parbox[c][2.5cm][c]{0.85\linewidth}{%
        \centering\footnotesize\textcolor{midgray}{[Missing image]}}}%
    }%
  }%
}

% --- Section count for navigation (must match actual \section{} count) ---
\setcounter{mlsystotalsections}{9}

\title{Performance Engineering}
\author{Vijay Janapa Reddi}
\institute{Harvard University}
\date{}

\begin{document}

% =============================================================================
% TITLE SLIDE
% =============================================================================
\mlsystitle{Performance Engineering}{Match the Software to the Silicon}{cover_optimization.png}

% =============================================================================
% LEARNING OBJECTIVES
% =============================================================================
\begin{frame}{Learning Objectives}
\note{[2 min] Walk through objectives. Emphasize that this chapter is about
closing the gap between theoretical hardware peak and achieved throughput.
Ask: ``Has anyone profiled a GPU workload and found it was only using 5\%
of compute?''}

\footnotesize
\begin{enumerate}\setlength\itemsep{0pt}
  \item Analyze the \textbf{Memory Wall} using the \textbf{Roofline Model} and diagnose compute-bound vs.\ memory-bound workloads
  \item Explain how \textbf{Operator Fusion} and \textbf{Tiling} (FlashAttention) overcome HBM bandwidth limits
  \item Apply \textbf{Precision Engineering} (FP8, INT4, block-wise quantization) to double effective bandwidth
  \item Compare \textbf{Graph Compilation} frameworks (torch.compile, XLA, TensorRT)
  \item Evaluate \textbf{Speculative Decoding} and \textbf{MoE} as algorithmic innovations
  \item Design an \textbf{optimization sequence} using the profile-optimize-reprofile loop
  \item Use \textbf{System Profiling} tools to identify the dominant bottleneck
\end{enumerate}

\end{frame}

% =============================================================================
\section{Memory Wall}
% =============================================================================

\begin{frame}{The Memory Wall}
\note{[3 min] Core insight: an H100 sits 95\% idle during LLM decode because
it starves for data. The Memory Wall is physics, not engineering. Every
technique in this chapter attacks the bytes-moved-to-ops-performed ratio.
Ask: ``If the GPU has 989 TFLOPS, why does LLM decode use only 2\%?''}

% --- Full-width diagram ---
\centering
\safeimg[width=0.92\textwidth,height=4.5cm,keepaspectratio]{memory-wall.pdf}

\vspace{0.15cm}
\mlsysalert{Key:}{Moving data costs 1,280$\times$ more energy from HBM than from SRAM. This is physics, not engineering.}

\end{frame}

\begin{frame}{The Widening Gap}
\note{[3 min] Each GPU generation increases compute faster than bandwidth.
The ridge point has grown 4x in 7 years. Operations near 200 FLOP/byte
were compute-bound on V100 but memory-bound on B200.
Ask: ``What does this mean for your optimization investments?''}

\small
\renewcommand{\arraystretch}{1.15}
{\footnotesize
\begin{tabular}{@{}lrrrr@{}}
  \toprule
  \textbf{GPU} & \textbf{Year} & \textbf{FP16 (TFLOPS)} & \textbf{HBM BW (TB/s)} & \textbf{Ridge (FLOP/B)} \\
  \midrule
  V100 & 2017 & 125  & 0.9  & 139 \\
  A100 & 2020 & 312  & 2.0  & 156 \\
  H100 & 2022 & 989  & 3.35 & 295 \\
  B200 & 2024 & 4,500 & 8.0  & 563 \\
  \bottomrule
\end{tabular}
}

\vspace{0.2cm}
\begin{mlsyscard}{errorstroke}
\textbf{36$\times$ compute growth} vs.\ \textbf{8.9$\times$ bandwidth growth} in 7 years.\\
{\footnotesize More workloads become memory-bound with each generation. Fusion and precision engineering become \emph{more} important, not less.}
\end{mlsyscard}

\end{frame}

% =============================================================================
\section{Roofline Model}
% =============================================================================

\begin{frame}{The Roofline Model}
\note{[3 min] The roofline gives a quantitative framework. Two ceilings:
bandwidth (sloped) and compute (flat). The ridge point is where they meet.
Most transformer inference falls left of the ridge.
Common error: students think more FLOPS always helps.}

% --- Full-width diagram ---
\centering
\safeimg[width=0.92\textwidth,height=4.8cm,keepaspectratio]{roofline-model.pdf}

\vspace{0.1cm}
\mlsysconcept{Diagnosis:}{If $\AI < \ridge$, you are memory-bound. Faster compute yields zero improvement.}

\end{frame}

\begin{frame}{Where ML Workloads Fall}
\note{[2 min] Three orders of magnitude separate LLM decode from batched GEMM.
The key insight is that the SAME model has different bottlenecks depending
on the operation and batch size.}

\footnotesize
\renewcommand{\arraystretch}{1.15}
\begin{tabular}{@{}lrll@{}}
  \toprule
  \textbf{Operation} & \textbf{AI (FLOP/B)} & \textbf{Regime} & \textbf{Optimization} \\
  \midrule
  LLM Decode (batch=1) & $\approx$2 & \textcolor{computestroke}{Deeply mem-bound} & Precision, fusion \\
  LayerNorm / GELU & $\approx$5 & \textcolor{computestroke}{Memory-bound} & Element-wise fusion \\
  Attention (naive) & $\approx$10 & \textcolor{computestroke}{Memory-bound} & FlashAttention \\
  GEMM $4096^2$ FP16 & $\approx$1,365 & \textcolor{errorstroke}{Compute-bound} & Tensor Cores \\
  Batched GEMM & $>$500 & \textcolor{errorstroke}{Compute-bound} & FP8 compute \\
  \bottomrule
\end{tabular}

\vspace{0.15cm}
\begin{mlsyscard}{crimson}
\textbf{Most transformer inference is memory-bound.} The H100 ridge point is 295 FLOP/byte at FP16. Only large GEMMs exceed it.
\end{mlsyscard}

\end{frame}

% --- ACTIVE LEARNING 1: Predict ---
\begin{frame}{Predict: Which Bottleneck Dominates?}
\note{[2 min] Give students 60 seconds. The answer is memory-bound:
batch-1 decode has AI $\approx$ 2, far below the ridge of 295.
Upgrading to a faster GPU (more TFLOPS) would yield zero improvement.}

\centering
\vspace{1.0cm}
{\Large\bfseries Think--Write--Share}

\vspace{0.5cm}
{\large A team upgrades from A100 to H100, expecting 3$\times$ faster\\
LLM inference at batch size 1. They see only 15\% improvement.}

\vspace{0.4cm}
{\normalsize Why? Which Iron Law term dominates?}

\vspace{0.4cm}
{\small\textcolor{midgray}{(60 seconds --- write your diagnosis before I reveal)}}

\end{frame}

% =============================================================================
\section{Operator Fusion}
% =============================================================================

\begin{frame}{Operator Fusion: Before and After}
\note{[3 min] Walk through left (unfused) vs right (fused). The key is
intermediate tensors Z1 and Z2 that the naive path writes to HBM and reads
back. The fused kernel keeps everything in SRAM.
Quantify: 32 MB of wasted traffic per layer.}

% --- Full-width diagram ---
\centering
\safeimg[width=0.92\textwidth,height=4.8cm,keepaspectratio]{operator-fusion.pdf}

\vspace{0.1cm}
\mlsysinsight{Impact:}{Fusion reduces HBM transfers from 5 to 2 per layer, saving 32 MB of redundant traffic.}

\end{frame}

\begin{frame}{Fusion Categories}
\note{[2 min] Three levels of increasing difficulty and impact.
Element-wise fusion is automatic. Reduction fusion combines normalization.
Operator-specific fusion (FlashAttention) is the big win but requires
custom kernel engineering. Cumulative effect: 60-80\% less HBM traffic.}

\scriptsize
\renewcommand{\arraystretch}{1.1}
\begin{tabular}{@{}llrl@{}}
  \toprule
  \textbf{Category} & \textbf{Example} & \textbf{Savings} & \textbf{Difficulty} \\
  \midrule
  \textbf{Element-wise} & Bias + GELU + Dropout & 64 MB/layer & Automatic \\
  \textbf{Reduction} & LayerNorm (mean + var) & 48 MB/layer & Moderate \\
  \textbf{Operator-specific} & FlashAttention & 4 GB/layer & Custom kernel \\
  \bottomrule
\end{tabular}

\vspace{0.1cm}
\begin{columns}[T]
  \begin{column}{0.48\textwidth}
    \begin{mlsyscard}{computestroke}
    {\scriptsize\textbf{CUDA Graphs} eliminate launch overhead:\\
    30+ kernels $\times$ 70 layers = 15 ms tax $\to$ $<$0.1 ms.}
    \end{mlsyscard}
  \end{column}
  \begin{column}{0.48\textwidth}
    \begin{mlsyscard}{datastroke}
    {\scriptsize\textbf{Combined:} Fusion + CUDA Graphs reduce HBM traffic 60--80\% and eliminate non-compute overhead.}
    \end{mlsyscard}
  \end{column}
\end{columns}

\end{frame}

\begin{frame}{FlashAttention: Tiled Attention}
\note{[3 min] The key insight: standard attention materializes an NxN matrix
in HBM (128 MB per head at seq=8192). FlashAttention tiles the computation
so this matrix never leaves SRAM. Online softmax makes this exact, not
approximate. Result: 65x reduction in HBM traffic.
Ask: ``Why can we not just use more HBM?''}

% --- Full-width diagram ---
\centering
\safeimg[width=0.92\textwidth,height=4.2cm,keepaspectratio]{flashattention.pdf}

\vspace{0.05cm}
{\footnotesize\textbf{Key Innovation:} Online softmax computes the \emph{exact} result tile-by-tile, never materializing $N \times N$.}

\end{frame}

% =============================================================================
\section{Precision}
% =============================================================================

\begin{frame}{Precision Engineering}
\note{[3 min] Halving precision doubles effective bandwidth -- the cheapest
optimization available. Walk through FP16 vs FP8 vs INT4 bit layouts.
The challenge: LLMs have outlier features that break uniform quantization.
Four solutions with different trade-offs.}

% --- Full-width diagram ---
\centering
\safeimg[width=0.92\textwidth,height=4.8cm,keepaspectratio]{precision-engineering.pdf}

\vspace{0.1cm}
\mlsysconcept{Principle:}{Fewer bytes per weight = more weights per HBM trip = higher effective bandwidth.}

\end{frame}

\begin{frame}{Quantization Techniques for LLMs}
\note{[3 min] Four approaches to the outlier problem. LLM.int8() handles
outliers at runtime. GPTQ uses Hessian for INT4. AWQ is fastest (minutes
vs hours). SmoothQuant enables W8A8. AWQ is the production default.
Common error: students think quantization always means quality loss.}

\footnotesize
\renewcommand{\arraystretch}{1.1}
\begin{tabular}{@{}lllll@{}}
  \toprule
  \textbf{Method} & \textbf{Precision} & \textbf{Approach} & \textbf{Calibration} & \textbf{Use Case} \\
  \midrule
  LLM.int8() & W8A-FP16 & Mixed-precision decomp. & Runtime & Safe default \\
  GPTQ & W4A-FP16 & Hessian-based & Hours & Max compression \\
  \textbf{AWQ} & W4A-FP16 & Activation-aware & \textbf{Minutes} & \textbf{Production default} \\
  SmoothQuant & W8A8 & Smooth activations & Offline & Compute-bound \\
  \bottomrule
\end{tabular}

\vspace{0.15cm}
\begin{columns}[T]
  \begin{column}{0.48\textwidth}
    \begin{mlsyscard}{routingstroke}
    \textbf{PTQ vs.\ QAT:}\\
    {\scriptsize PTQ first (4 GPU-hours for 70B).\\
    QAT only if quality fails (3 days on 8$\times$H100).\\
    QLoRA bridges the gap.}
    \end{mlsyscard}
  \end{column}
  \begin{column}{0.48\textwidth}
    \begin{mlsyscard}{errorstroke}
    \textbf{Weight-only vs.\ W+A:}\\
    {\scriptsize INT4 weight-only at small batch (latency).\\
    FP8 W+A at large batch (throughput).\\
    Many systems use both.}
    \end{mlsyscard}
  \end{column}
\end{columns}

\end{frame}

% =============================================================================
\section{Compilation}
% =============================================================================

\begin{frame}{Graph Compilation Pipeline}
\note{[3 min] Five stages from Python to optimized kernels. The key insight:
compilers automate known optimizations. They cannot discover new algorithms
like FlashAttention or speculative decoding.
Walk through the pipeline left-to-right.}

% --- Full-width diagram ---
\centering
\safeimg[width=0.92\textwidth,height=4.5cm,keepaspectratio]{graph-compilation.pdf}

\vspace{0.1cm}
\mlsysalert{Limitation:}{Compilers automate \emph{known} optimizations; human engineers discover \emph{new} ones.}

\end{frame}

\begin{frame}{Three Compilers, Three Contexts}
\note{[2 min] torch.compile for development (seconds, flexible), XLA for
large-scale training (minutes, whole-program), TensorRT for production
inference (30-60 min, maximum throughput). The choice depends on deployment.
Common error: using TensorRT for research iteration.}

\footnotesize
\renewcommand{\arraystretch}{1.15}
\begin{tabular}{@{}llll@{}}
  \toprule
  & \textbf{torch.compile} & \textbf{XLA} & \textbf{TensorRT} \\
  \midrule
  \textbf{Ecosystem} & PyTorch + GPU & JAX + TPU & NVIDIA inference \\
  \textbf{Compile time} & Seconds & Minutes & 30--60 min \\
  \textbf{Speedup} & 10--40\% & 55--65\% MFU & 1.3--1.8$\times$ \\
  \textbf{Flexibility} & Graph breaks OK & Static graph only & Static shapes \\
  \textbf{Best for} & Dev \& research & Training at scale & Prod serving \\
  \bottomrule
\end{tabular}

\vspace{0.15cm}
\mlsysconcept{Trade-off:}{Flexibility $\leftrightarrow$ Performance. More constraints $=$ more optimization opportunity.}

\end{frame}

% --- ACTIVE LEARNING 2: Exercise ---
\begin{frame}{Your Turn: Diagnose and Optimize}
\note{[3 min] Give 90 seconds. Answer: Memory-bound (AI=2, ridge=295).
Step 1: INT4 quantization (4x BW). Step 2: FlashAttention.
Step 3: torch.compile. Upgrading FLOPS is useless.}

\footnotesize
\begin{columns}[T]
  \begin{column}{0.55\textwidth}
    {\small\bfseries Diagnose the Bottleneck}

    \vspace{0.1cm}
    A 70B LLM serves batch-1 decode on an H100:
    \begin{itemize}\setlength\itemsep{0pt}
      \item Arithmetic Intensity $\approx$ 2 FLOP/byte
      \item H100 ridge point = 295 FLOP/byte
      \item Achieved BW = 2.8 TB/s (84\% of peak)
      \item MFU = 2\%
    \end{itemize}

    \vspace{0.05cm}
    \textbf{1.} Is this compute-bound or memory-bound?\\
    \textbf{2.} Name the top 3 optimizations \emph{in order}.\\
    \textbf{3.} Would upgrading to B200 (2$\times$ FLOPS) help?

    {\scriptsize\textcolor{midgray}{(90 seconds --- compare with a neighbor)}}
  \end{column}
  \begin{column}{0.42\textwidth}
    \pause
    \begin{mlsyscard}{datastroke}
    {\scriptsize\textbf{Solution:}\\[0.05cm]
    1. \textcolor{computestroke}{\textbf{Memory-bound}} (AI $\ll$ ridge)\\[0.05cm]
    2. a) INT4 quant (4$\times$ BW)\\
    \quad b) FlashAttention (65$\times$)\\
    \quad c) torch.compile (10--40\%)\\[0.05cm]
    3. B200 has $2\times$ FLOPS but workload is BW-limited. \textbf{No help.}
    }
    \end{mlsyscard}
  \end{column}
\end{columns}

\end{frame}

% =============================================================================
\section{Algorithmic}
% =============================================================================

\begin{frame}{Speculative Decoding}
\note{[3 min] The key insight: LLM decode is memory-bound because each token
loads the full weight matrix. Speculative decoding uses a small draft model
to predict k tokens, then the target model verifies all k in one pass.
The target's weights are loaded once for k tokens instead of once per token.
Speedup: 1.5-3x. Best at batch=1.}

\footnotesize
\begin{columns}[T]
  \begin{column}{0.55\textwidth}
    \textbf{Problem:} Each decode step loads full weights from HBM for \emph{one} token. AI $\approx$ 2 FLOP/byte.

    \vspace{0.1cm}
    \textbf{Speculative Decoding:}
    \begin{enumerate}\setlength\itemsep{0pt}
      \item \textcolor{datastroke}{Draft model} generates $k$ candidates (fast)
      \item \textcolor{computestroke}{Target model} verifies all $k$ in \textbf{one} pass
      \item Accepted tokens: mathematically identical output
      \item Rejected tokens: ``wasted'' compute (but compute is free!)
    \end{enumerate}

    \vspace{0.05cm}
    \mlsysconcept{Trade:}{Abundant compute for scarce bandwidth.}
  \end{column}
  \begin{column}{0.42\textwidth}
    \begin{mlsyscard}{computestroke}
    {\scriptsize\textbf{Why it works:} Target loads weights \textbf{once} for $k$ tokens instead of once per token. Effective AI increases by $k\times$. Speedup: \textbf{1.5--3$\times$} typical.}
    \end{mlsyscard}

    \vspace{0.1cm}
    \begin{mlsyscard}{errorstroke}
    {\scriptsize\textbf{Caveat:} Best at batch=1 (memory-bound). At large batches, speculation adds overhead.}
    \end{mlsyscard}
  \end{column}
\end{columns}

\end{frame}

\begin{frame}{Mixture of Experts (MoE)}
\note{[3 min] MoE decouples model capacity from inference cost. A 1T dense
model would need 2 TB of memory. An MoE with 370B total but 37B active
per token achieves similar quality at 27x less compute per token.
DeepSeek-V3: 671B total, 37B active. Challenge: AllToAll communication
and load balancing.}

\footnotesize
\begin{columns}[T]
  \begin{column}{0.55\textwidth}
    \textbf{Break the capacity-cost link:}
    \begin{itemize}\setlength\itemsep{0pt}
      \item Replace FFN with $N$ parallel ``expert'' FFNs
      \item Lightweight \textbf{router} selects top-$k$ experts/token
      \item Only $k/N$ of parameters activate per token
    \end{itemize}

    \vspace{0.05cm}
    {\scriptsize
    \renewcommand{\arraystretch}{1.05}
    \begin{tabular}{@{}lrr@{}}
      \toprule
      & \textbf{Dense 1T} & \textbf{MoE 370B/37B} \\
      \midrule
      Total params  & 1,000B & 370B \\
      Active/token  & 1,000B & 37B \\
      Memory (FP16) & 2,000 GB & 740 GB \\
      Compute ratio & 1$\times$ & 27$\times$ less \\
      \bottomrule
    \end{tabular}
    }
  \end{column}
  \begin{column}{0.42\textwidth}
    \begin{mlsyscard}{datastroke}
    {\scriptsize\textbf{DeepSeek-V3:} 671B total, 37B active. 256 experts, top-8 routing. Near GPT-4 quality at fraction of cost.}
    \end{mlsyscard}

    \vspace{0.08cm}
    \begin{mlsyscard}{errorstroke}
    {\scriptsize\textbf{System Challenges:}\\
    $\bullet$ AllToAll communication\\
    $\bullet$ Load balancing across experts\\
    $\bullet$ Training instability ($>$64 experts)
    }
    \end{mlsyscard}
  \end{column}
\end{columns}

\end{frame}

% =============================================================================
\section{Profiling}
% =============================================================================

\begin{frame}{System Profiling: Measure First}
\note{[3 min] The engineer who spent 2 weeks on a CUDA kernel that yielded
zero improvement because the system was I/O bound. Profiling has four levels.
The key workflow: app-level to find symptoms, trace-level to find gaps,
kernel-level to diagnose. Never skip levels.}

\footnotesize
\textbf{Four-Level Profiling Hierarchy}

\vspace{0.05cm}
\scriptsize
\renewcommand{\arraystretch}{1.05}
\begin{tabular}{@{}llp{5cm}@{}}
  \toprule
  \textbf{Level} & \textbf{Tool} & \textbf{What It Reveals} \\
  \midrule
  Application & Custom metrics & Tokens/s, TTFT, P99 latency \\
  Distributed & Nsight Systems & Comm patterns, load imbalance \\
  Trace & PyTorch Profiler & Kernel timeline, GPU idle gaps \\
  Kernel & Nsight Compute & Achieved BW, occupancy, cache misses \\
  \bottomrule
\end{tabular}

\vspace{0.05cm}
\scriptsize
\begin{columns}[T]
  \begin{column}{0.48\textwidth}
    \begin{mlsyscard}{crimson}
    {\scriptsize\textbf{Key Metrics:}\\
    $\bullet$ \textbf{MFU}: training (40--60\% typical)\\
    $\bullet$ \textbf{MBU}: inference (85\% = near-optimal)\\
    $\bullet$ \textbf{TTFT}: time-to-first-token ($<$500 ms)\\
    $\bullet$ \textbf{ITL}: inter-token latency ($<$50 ms)}
    \end{mlsyscard}
  \end{column}
  \begin{column}{0.48\textwidth}
    \begin{mlsyscard}{errorstroke}
    {\scriptsize\textbf{MFU misleads for inference.} Decode may show 2\% MFU but 85\% MBU --- the system is near-optimal. Use MBU for inference, MFU for training.}
    \end{mlsyscard}
  \end{column}
\end{columns}

\end{frame}

\begin{frame}{Iron Law Diagnostic Flowchart}
\note{[3 min] Walk through the flowchart: profile first, classify the
bottleneck, then apply the right technique. The critical insight: each
optimization shifts the bottleneck, requiring reprofiling.
Common error: applying all optimizations without measuring.}

% --- Full-width diagram ---
\centering
\safeimg[width=0.92\textwidth,height=4.5cm,keepaspectratio]{diagnostic-flow.pdf}

\vspace{0.1cm}
\mlsysalert{Key:}{Every optimization shifts the bottleneck. The profile-optimize-reprofile loop never ends.}

\end{frame}

% --- ACTIVE LEARNING 3: Discussion ---
\begin{frame}{Discussion: The Optimization Trap}
\note{[3 min] Turn-and-talk. Students discuss in pairs for 90 seconds.
The point: you must reprofile after each step. INT4 may shift the
bottleneck from memory to compute, making further bandwidth optimization
useless.}

\centering
\vspace{0.8cm}
{\large\bfseries Turn and Talk \textcolor{midgray}{(90 seconds)}}

\vspace{0.5cm}
{\large A team applies INT4 quantization (4$\times$ bandwidth savings)\\
to their memory-bound serving system.\\[0.2cm]
Throughput improves by only 2$\times$, not 4$\times$.}

\vspace{0.4cm}
\alert{What happened? What should they do next?}

\vspace{0.3cm}
{\small\textcolor{midgray}{Hint: think about what the bottleneck became \emph{after} quantization.}}

\end{frame}

% =============================================================================
\section{The Playbook}
% =============================================================================

\begin{frame}{The Optimization Playbook}
\note{[3 min] Walk through the five-step sequence. The principle: passive
optimizations first (compiler, library swaps), active ones last
(algorithmic changes). Each step validated by reprofiling.
The order is critical: torch.compile costs 1 line, speculation costs weeks.}

% --- Full-width diagram ---
\centering
\safeimg[width=0.92\textwidth,height=4.5cm,keepaspectratio]{optimization-playbook.pdf}

\vspace{0.1cm}
\mlsysconcept{Principle:}{Passive optimizations first $\to$ active changes last. Reprofile after every step.}

\end{frame}

\begin{frame}{Case Study: 70B LLM Optimization}
\note{[3 min] Walk through the compounding effects. INT4 frees memory for
larger batches, which changes arithmetic intensity, which determines the
next optimization. Emphasize the iterative nature.
If short: show just the summary table.}

\footnotesize
\renewcommand{\arraystretch}{1.1}
\begin{tabular}{@{}lrrrl@{}}
  \toprule
  \textbf{Step} & \textbf{Throughput} & \textbf{Latency} & \textbf{Gain} & \textbf{Mechanism} \\
  \midrule
  Baseline (FP16) & 80 tok/s & 50 ms ITL & --- & Unoptimized \\
  + torch.compile & 100 tok/s & 42 ms & 1.25$\times$ & Fusion + kernel select \\
  + AWQ INT4 & 280 tok/s & 35 ms & 2.8$\times$ & 4$\times$ less weight BW \\
  + FlashAttention & 350 tok/s & 30 ms & 1.25$\times$ & O(N) attention HBM \\
  + KV INT8 cache & 600 tok/s & 45 ms & 1.7$\times$ & 2$\times$ batch size \\
  + Spec.\ decoding & 600 tok/s & 22 ms & lat.\ 2$\times$ & Draft model parallel \\
  \bottomrule
\end{tabular}

\vspace{0.15cm}
\begin{mlsyscard}{crimson}
\textbf{7.5$\times$ total throughput}, \textbf{2.3$\times$ latency reduction} from the same hardware.\\
{\scriptsize Each optimization shifted the bottleneck. INT4 $\to$ larger batch $\to$ compute-bound $\to$ new optimization target.}
\end{mlsyscard}

\end{frame}

% =============================================================================
\section{Wrap-Up}
% =============================================================================

\begin{frame}{Fallacies}
\note{[2 min] Three common misconceptions with quantitative evidence.}

\small
\textbf{Fallacy:} \textit{More FLOPS means faster inference.}\\
{\footnotesize Most inference is memory-bound. A GPU with 2$\times$ FLOPS but the same bandwidth generates tokens at the same speed. Correct metric: bandwidth, not FLOPS.}

\vspace{0.15cm}
\textbf{Fallacy:} \textit{FP8 always halves training time.}\\
{\footnotesize FP8 doubles peak TFLOPS and effective BW, but only for precision-sensitive operations. Element-wise ops are launch-overhead bound. Communication-bound steps gain nothing.}

\vspace{0.15cm}
\textbf{Fallacy:} \textit{Graph compilers eliminate manual kernel engineering.}\\
{\footnotesize FlashAttention required recognizing that attention could be tiled with online softmax --- an insight beyond any compiler's rewrite rules. Compilers automate known optimizations; humans discover new ones.}

\end{frame}

\begin{frame}{Pitfalls}
\note{[2 min] Four operational pitfalls. Focus on the first two if short on time.}

\scriptsize
\textbf{Pitfall:} \textit{Optimizing the largest kernel while ignoring the long tail.}\\
When the GEMM is near-optimal, dozens of smaller ops can consume $>$50\% of time. Graph compilation addresses this tail.

\vspace{0.08cm}
\textbf{Pitfall:} \textit{Quantizing everything to the lowest precision.}\\
INT4 weights + INT4 KV + FP8 activations can degrade quality invisibly. Targeted quantization: aggressive for KV cache, conservative for first/last layers.

\vspace{0.08cm}
\textbf{Pitfall:} \textit{Measuring throughput without measuring quality.}\\
200 tok/s with 15\% quality loss $\neq$ better than 100 tok/s at full quality. Optimize on the throughput vs.\ quality Pareto frontier.

\vspace{0.08cm}
\textbf{Pitfall:} \textit{A single profiling run characterizes performance.}\\
Thermal throttling (10--15\%), fragmentation, and congestion vary. Profile under sustained conditions.

\end{frame}

% --- RETRIEVAL PRACTICE ---
\begin{frame}{What Were the Key Ideas?}
\note{[2 min] Retrieval practice. Students write 90 seconds, no notes.
Walk around the room. Do NOT show the next slide yet.}

\centering
\vspace{1.5cm}
{\Large\bfseries Close your notes.}

\vspace{0.8cm}
{\large Write down the \textbf{4 most important concepts} from today.}

\vspace{0.8cm}
{\normalsize\textcolor{midgray}{90 seconds --- no peeking.}}

\end{frame}

\begin{frame}{Key Takeaways}
\note{[2 min] Reveal. Walk through each bullet. Quantitative anchors:
295 FLOP/byte ridge, 65x FlashAttention, 4x INT4, 7.5x total in case study.}

\scriptsize
\begin{itemize}\setlength\itemsep{0pt}
  \item \textbf{Memory Wall}: Most transformer inference is memory-bound (AI $\ll$ 295 FLOP/byte). Reduce bytes moved, not operations computed.
  \item \textbf{Roofline Model}: Diagnose compute-bound vs.\ memory-bound \emph{before} optimizing. Wrong-term optimization yields zero improvement.
  \item \textbf{Operator Fusion}: FlashAttention reduces attention HBM traffic 65$\times$ via tiling + online softmax. CUDA Graphs eliminate launch overhead.
  \item \textbf{Precision}: INT4 (AWQ) gives 4$\times$ effective bandwidth. Block-wise quantization handles LLM outlier features.
  \item \textbf{Graph Compilation}: torch.compile (1 line, 10--40\%), XLA (55--65\% MFU), TensorRT (1.3--1.8$\times$).
  \item \textbf{Algorithmic}: Speculative decoding trades compute for latency (1.5--3$\times$). MoE decouples capacity from cost.
  \item \textbf{Profile-Optimize-Reprofile}: Every optimization shifts the bottleneck. The iterative loop is the discipline.
\end{itemize}

\end{frame}

\begin{frame}{References}
\note{[1 min] Point students to canonical papers.}

\small
\mlsysref{Dao+22}{Dao et al. ``FlashAttention: Fast and Memory-Efficient Exact Attention.'' NeurIPS 2022.}
\mlsysref{Williams+09}{Williams, Waterman, Patterson. ``Roofline: An Insightful Visual Performance Model.'' CACM 2009.}
\mlsysref{Dettmers+22}{Dettmers et al. ``LLM.int8(): 8-bit Matrix Multiplication for Transformers at Scale.'' NeurIPS 2022.}
\mlsysref{Lin+24}{Lin et al. ``AWQ: Activation-Aware Weight Quantization.'' MLSys 2024.}
\mlsysref{Leviathan+23}{Leviathan et al. ``Fast Inference from Transformers via Speculative Decoding.'' ICML 2023.}
\mlsysref{DeepSeek24}{DeepSeek-AI. ``DeepSeek-V3 Technical Report.'' 2024.}

\end{frame}

\begin{frame}{Next Lecture: Edge Intelligence}
\note{[1 min] Forward hook. Performance engineering optimizes within the
datacenter. But what happens when the model must run on devices with
1,000x less memory and 10,000x less compute? Edge Intelligence brings
these constraints to the extreme.}

\footnotesize
\begin{columns}[c]
  \begin{column}{0.30\textwidth}
    \centering
    \begin{mlsyscard}{computestroke}
    {\large\bfseries Cloud}\\[0.1cm]
    {\footnotesize 80 GB HBM\\989 TFLOPS\\700 W}
    \end{mlsyscard}
  \end{column}
  \begin{column}{0.30\textwidth}
    \centering
    \begin{mlsyscard}{routingstroke}
    {\large\bfseries Edge}\\[0.1cm]
    {\footnotesize 32 GB LPDDR5\\100 TOPS\\15--40 W}
    \end{mlsyscard}
  \end{column}
  \begin{column}{0.30\textwidth}
    \centering
    \begin{mlsyscard}{errorstroke}
    {\large\bfseries TinyML}\\[0.1cm]
    {\footnotesize 384 KB SRAM\\0.5 GFLOPS\\100 mW}
    \end{mlsyscard}
  \end{column}
\end{columns}

\vspace{0.3cm}
\centering
{\normalsize When the model \emph{must} run where the data lives,}\\[0.1cm]
{\small\textbf{every optimization technique from today becomes a survival skill.}}

\end{frame}

\end{document}
